\chapter{यांत्रिक दोलित्र और समय का मापन}\label{ch:g12:oscillators-and-time}

दादाजी की घड़ी टिकटिक करती है: हर टिक पीतल के एक लोलक का \omterm{def:g10:universal-gravitation:weight}{ऊर्ध्वाधर} पार करना है,
जो नए साल से अब तक डेढ़ करोड़ बार दोहराया जा चुका है। आपकी कलाई-घड़ी में क्वार्ट्ज़
की एक पतली पर्त \omterm{def:g10:orders-of-magnitude:unit}{सेकंड} में \num{32768} बार लचकती है; और सीज़ियम का एक \omterm{def:g11:fundamental-interactions:ladder}{परमाणु}
उससे नौ अरब गुना तेज़ गुनगुनाता है। जो कुछ झूलता है वह समय रख सकता है: यह
अध्याय उसी झूलने के पीछे चलता है, हाइगेंस से लेकर उपग्रहों तक।

\section{मुक्त दोलन}

\begin{definition}[दोलित्र, आवर्तकाल, आवृत्ति, आयाम]\label{def:g12:oscillators-and-time:oscillator}
\emph{यांत्रिक दोलित्र}\index{दोलित्र} वह निकाय है जो किसी स्थायी
\emph{साम्यावस्था स्थिति}\index{साम्यावस्था स्थिति} से हटाकर छोड़ देने पर उसके
चारों ओर आगे-पीछे झूलता रहे: ये \emph{मुक्त दोलन}\index{मुक्त दोलन} हैं। गति
\emph{आवर्तकाल}\index{आवर्तकाल} $T_0$ (\unit{s}) के बाद दोहरा जाती है;
\emph{आवृत्ति}\index{आवृत्ति} $f_0 = 1/T_0$ (\unit{Hz}, \cref{ch:g12:mechanical-waves}); और
\emph{आयाम}\index{आयाम} अधिकतम विस्थापन है।
\end{definition}

\begin{example}[दोलित्रों का चिड़ियाघर]\label{ex:g12:oscillators-and-time:zoo}
खेल के मैदान का झूला: $T_0 \approx \qty{3}{s}$। गिटार का A तार: $f_0 = \qty{110}{Hz}$। क्वार्ट्ज़ का क्रिस्टल:
\qty{32768}{Hz}। \omterm{def:g10:orders-of-magnitude:unit}{सेकंड} को परिभाषित करता सीज़ियम का दोलन: लगभग \qty{9.2e9}{Hz}। दस कोटियाँ, और
काम एक ही: गिनना।
\end{example}

\section{सरल लोलक}

\begin{definition}[सरल लोलक]\label{def:g12:oscillators-and-time:pendulum}
\emph{सरल लोलक}\index{सरल लोलक}: लंबाई $L$ के अवितान्य और नगण्य द्रव्यमान
वाले तार पर लटका द्रव्यमान $m$ का छोटा गोलक, जो किसी \omterm{def:g10:universal-gravitation:weight}{ऊर्ध्वाधर} तल में
झूलता हो। गोलक पर लगते दोनों बलों (\cref{ch:g12:newtons-laws}) में से \omterm{def:g10:inertia:inventory}{तनाव} गति के
लंबवत है, जबकि \omterm{def:g10:relative-motion:trajectory}{गतिपथ} के अनुदिश भार का घटक सबसे निचले बिंदु की
ओर पीछे संकेत करता है: यानी भार ही \emph{प्रत्यानयन
बल}% \index{प्रत्यानयन बल} है, जो विस्थापन का विरोध करता है।
\end{definition}

\begin{omfigure}
\begin{tikzpicture}[scale=0.8, font=\small]
  \coordinate (O) at (0,0); \draw (-0.6,0) -- (0.6,0);
  \foreach \x in {-0.5,-0.25,0,0.25,0.5} \draw (\x,0) -- ++(0.14,0.14);
  \fill (O) circle (1.5pt); \draw[dashed] (O) -- (0,-3.4);
  \draw[dashed, omExo] (-120:3.0) arc (-120:-60:3.0);
  \draw (O) -- (-60:3.0) node[midway, above right] {$L$};
  \draw (-90:1.0) arc (-90:-60:1.0); \node at (-75:1.32) {$\theta$};
  \coordinate (B) at (1.5,-2.598); \fill[omDef] (B) circle (2.5pt);
  \draw[-Stealth, omDef, thick] (B) -- ++(-0.75,1.299) node[above left] {$\vect T$};
  \draw[-Stealth, thick] (B) -- ++(0,-1.8) node[below] {$\vect P$};
  \draw[-Stealth, omThm, thick] (B) -- ++(-0.779,-0.45) node[left] {$\vect{P_t}$};
  \draw[dashed, omExo] (0.721,-3.048) -- (1.5,-4.398)
    (B) -- ++(0.779,-1.35) (2.279,-3.948) -- (1.5,-4.398);
\end{tikzpicture}

{\small गोलक पर बल: \omterm{def:g10:inertia:inventory}{तनाव} $\vect T$ गति के लंबवत है; और
भार का स्पर्शरेखीय हिस्सा $\vect{P_t}$ सदा \omterm{def:g10:universal-gravitation:weight}{ऊर्ध्वाधर} की ओर पीछे संकेत करता
है --- यही \omterm{def:g12:oscillators-and-time:pendulum}{प्रत्यानयन बल} है।}
\end{omfigure}

\begin{proposition}[सरल लोलक का आवर्तकाल]\label{prop:g12:oscillators-and-time:pendulum-period}
\index{सरल लोलक!आवर्तकाल}
छोटे आयामों (लगभग $15^\circ$ से नीचे) के लिए \omterm{def:g12:oscillators-and-time:oscillator}{आवर्तकाल} है
\[
T_0 = 2\pi \sqrt{\frac{L}{g}}, \qquad g = \qty{9.81}{m/s^2}.
\]
वह न गोलक के द्रव्यमान पर निर्भर करता है, न --- जब तक कोण छोटा रहे ---
आयाम पर: छोटे झूले \emph{समकालिक}\index{समकालिकता} होते हैं, और सबमें
एक ही समय लगता है।
\end{proposition}
\admitted

\begin{remark}[विमीय विश्लेषण, और $2\pi$ कहाँ रहता है]\label{rem:g12:oscillators-and-time:dimensional}
सूत्र का अंदाज़ा लगाया जा सकता है। $m$ (\unit{kg}), $L$ (\unit{m}) और $g$
(\unit{m/s^2}) में से समय के मात्रक वाला एकमात्र संयोजन $\sqrt{L/g}$ है --- और
द्रव्यमान उसमें आ ही \emph{नहीं} सकता, क्योंकि उसे काटने के लिए और कोई आँकड़ा
\omterm{def:g10:orders-of-magnitude:unit}{किलोग्राम} ढोता ही नहीं। पर विमाएँ आगे लगती शुद्ध संख्या नहीं दे सकतीं:
$2\pi$ गति-समीकरण हल करने से आता है, जो वर्ष 1 के खंड में किया गया है। और
छोटे-कोण की शर्त भी छोड़ी नहीं जा सकती: $30^\circ$ पर सच्चा \omterm{def:g12:oscillators-and-time:oscillator}{आवर्तकाल}
$2\%$ लंबा हो जाता है, और समकालिकता टूट जाती है।
\end{remark}

\begin{example}[सेकंड वाला लोलक]\label{ex:g12:oscillators-and-time:seconds}
कौन-सी लंबाई \omterm{def:g10:orders-of-magnitude:unit}{सेकंड} की चाल देती है, $T_0 = \qty{2.00}{s}$, यानी हर ओर एक टिक? हल करने पर
$L = g\,T_0^2/4\pi^2 = 9.81 \times 2.00^2/4\pi^2 \approx
\qty{0.994}{m}$। यानी लगभग एक \omterm{def:g10:orders-of-magnitude:unit}{मीटर}: लोलक-घड़ियाँ इसीलिए लंबी होती हैं कि \omterm{def:g10:orders-of-magnitude:unit}{सेकंड} लंबा
होता है।
\end{example}

\section{स्प्रिंग--द्रव्यमान दोलित्र}

\begin{definition}[स्प्रिंग का प्रत्यानयन बल]\label{def:g12:oscillators-and-time:spring}
द्रव्यमान $m$ का एक फिसलक क्षैतिज घर्षणहीन पटरी पर सरकता है और
\emph{दृढ़ता}\index{दृढ़ता} $k$ (\unit{N/m}) की स्प्रिंग से बँधा है।
साम्यावस्था से उसका विस्थापन $x$ लीजिए। खिंची या दबी हुई स्प्रिंग फिसलक को
बल $F = -kx$ से वापस खींचती या धकेलती है: यानी विस्थापन के समानुपाती और
उसके विरुद्ध --- फिर वही \omterm{def:g12:oscillators-and-time:pendulum}{प्रत्यानयन बल}, बस अब प्रत्यास्थता से मिला
हुआ।
\end{definition}

\begin{proposition}[गति-समीकरण और उसका हल]\label{prop:g12:oscillators-and-time:spring-period}
\index{स्प्रिंग--द्रव्यमान दोलित्र}
पटरी के अनुदिश \omterm{thm:g12:newtons-laws:second}{न्यूटन का दूसरा नियम} (\cref{ch:g12:newtons-laws}) $m \frac{d^2x}{dt^2} = -kx$ देता है, यानी
$\frac{d^2x}{dt^2} =
-\frac{k}{m}\,x$। अचर $A$ और $\varphi$ चाहे जो हों,
\[
x(t) = A \cos\!\left(\frac{2\pi t}{T_0} + \varphi\right),
\qquad T_0 = 2\pi\sqrt{\frac{m}{k}},
\]
इस समीकरण को संतुष्ट करता है: यानी \omterm{def:g12:oscillators-and-time:oscillator}{आवर्तकाल} $T_0$ का दोलन, जो
आयाम से स्वतंत्र है।
\end{proposition}

\begin{proof}
अवकलन करने पर $\frac{dx}{dt} = -\frac{2\pi}{T_0}A \sin(\frac{2\pi
t}{T_0} + \varphi)$; और फिर $\frac{d^2x}{dt^2} = -(\frac{2\pi}{T_0})^2
A\cos(\frac{2\pi t}{T_0} + \varphi) = -(\frac{2\pi}{T_0})^2\,x(t)$, जो ठीक तभी $-(k/m)\,x$ है जब $(2\pi/T_0)^2 = k/m$।
\end{proof}

\begin{remark}[गतियाँ जितनी हैं, सब यही हैं]\label{rem:g12:oscillators-and-time:uniqueness}
फिसलक की \emph{हर} गति इन्हीं फलनों में से एक है, यह मान लिया गया है और वर्ष 1
के खंड में सिद्ध किया गया है। और $2\pi\sqrt{m/k}$ में $g$ कहीं नहीं है: कड़ी स्प्रिंग
तेज़ झूलती है, भारी द्रव्यमान धीमा, और चंद्रमा पर भी वही।
\end{remark}

\begin{definition}[आयाम और कला]\label{def:g12:oscillators-and-time:phase}
$x(t) = A\cos(2\pi t/T_0 + \varphi)$ में धनात्मक अचर $A$ आयाम है --- छोर $x = \pm A$ पर हैं --- और
$2\pi t/T_0 +
\varphi$ \emph{कला}\index{कला} है, जो हर \omterm{def:g12:oscillators-and-time:oscillator}{आवर्तकाल} में $2\pi$ बढ़ जाती
है; \emph{आरंभिक कला}\index{आरंभिक कला} $\varphi$ शुरुआत दर्ज करती है: $x = A$ से
\omterm{def:g10:relative-motion:frame}{विराम} में छोड़ा जाए तो $\varphi = 0$। जैसे $k = \qty{20}{N/m}$ पर $m = \qty{0.20}{kg}$ $T_0 =
2\pi\sqrt{0.20/20} = \qty{0.63}{s}$ देता है; और
\qty{5.0}{cm} तक खींचकर छोड़ने पर सेंटीमीटर में $x(t) = 5.0\cos(2\pi t/0.63)$।
\end{definition}

\section{ऊर्जा, अवमंदन, अनुनाद}

\begin{proposition}[दोलित्र की ऊर्जा]\label{prop:g12:oscillators-and-time:energy}
$x$ जितनी खिंची स्प्रिंग प्रत्यास्थ स्थितिज ऊर्जा $E_p =
\tfrac12 kx^2$ संचित
रखती है (\cref{ch:g12:work-and-energy} में इसका औचित्य है)। घर्षणहीन दोलन के दौरान \omterm{def:g11:mechanical-energy:em}{यांत्रिक
ऊर्जा}
\[
E_m = E_k + E_p = \tfrac12 m v^2 + \tfrac12 k x^2 = \tfrac12 k A^2
\]
अचर रहती है: छोरों पर पूरी प्रत्यास्थ, साम्यावस्था पर पूरी गतिज। लोलक भी वही
युगल-गान $E_k$ और गुरुत्वीय $E_p$ के बीच गाता है (\cref{ch:g11:mechanical-energy})।
\end{proposition}

\begin{proof}
$v = \frac{dx}{dt}$ के साथ $E_m$ का अवकलज गति-समीकरण से $mv\frac{dv}{dt}
+ kx\frac{dx}{dt} = v\,(m\frac{d^2x}{dt^2} + kx) = 0$ हो जाता है; और मान
$\tfrac12 kA^2$ किसी छोर पर पढ़ लिया जाता है।
\end{proof}

\begin{omfigure}
\begin{tikzpicture}
\begin{axis}[omaxis, width=8.6cm, height=4.6cm, xlabel={$t/T_0$},
    ylabel={ऊर्जा $/\,E_m$}, xmin=0, xmax=1, ymin=0, ymax=1.3,
    legend pos=south east]
  \addplot[omDef, thick, domain=0:1, samples=120] {cos(360*x)^2};
  \addplot[omProp, thick, domain=0:1, samples=120] {sin(360*x)^2};
  \addplot[omThm, dashed, thick, domain=0:1, samples=2] {1};
  \legend{$E_p$, $E_k$, $E_m$}
\end{axis}
\end{tikzpicture}

{\small एक \omterm{def:g12:oscillators-and-time:oscillator}{आवर्तकाल}, $x = A$ से छोड़े जाने पर: ऊर्जा
प्रत्यास्थ से गतिज और वापस, दो बार छलकती है, जबकि योग $\tfrac12 kA^2$ पर टँका रहता
है।}
\end{omfigure}

\begin{definition}[अवमंदित दोलन]\label{def:g12:oscillators-and-time:damping}
\omterm{def:g10:inertia:inventory}{घर्षण} हर असली \omterm{def:g12:oscillators-and-time:oscillator}{दोलित्र} से \omterm{def:g11:mechanical-energy:em}{यांत्रिक ऊर्जा} खींच लेता
है: उसके \omterm{def:g12:oscillators-and-time:oscillator}{मुक्त दोलन} \emph{अवमंदित}\index{अवमंदन} होते हैं। हलके
अवमंदन में निकाय अब भी दोलन करता है, आयाम मरता जाता है पर दोहराव का समय लगभग
वही रहता है, यानी \emph{छद्म-आवर्तकाल}% \index{छद्म-आवर्तकाल},
जो $T_0$ के पास ही होता है; भारी अवमंदन (शहद) में वह बिना कभी छोर पार किए
रेंगते हुए साम्यावस्था पर लौट आता है; और आयाम के वर्ग के समानुपाती
ऊर्जा ऊष्मा में \omterm{def:g11:mechanical-energy:dissipation}{क्षय} हो जाती है।
\end{definition}

\begin{omfigure}
\begin{tikzpicture}
\begin{axis}[omaxis, width=8.6cm, height=4.6cm, xlabel={$t$ (s)},
    ylabel={$x$ (cm)}, xmin=0, xmax=10, ymin=-2.6, ymax=2.6,
    legend pos=south east]
  \addplot[omExo, thin, domain=0:10, samples=240] {2*cos(deg(pi*x))};
  \addplot[omDef, thick, domain=0:10, samples=240] {2*exp(-x/4)*cos(deg(pi*x))};
  \legend{undamped, damped}
  \addplot[omThm, dashed, domain=0:10, samples=60, forget plot] {2*exp(-x/4)};
  \addplot[omThm, dashed, domain=0:10, samples=60, forget plot] {-2*exp(-x/4)};
\end{axis}
\end{tikzpicture}

{\small \omterm{def:g12:oscillators-and-time:damping}{अवमंदित} और हलके \omterm{def:g12:oscillators-and-time:damping}{अवमंदित} दोलन: \omterm{def:g12:oscillators-and-time:oscillator}{आवर्तकाल} लगभग वही
(यहाँ \qty{2.0}{s}), और \omterm{def:g12:oscillators-and-time:damping}{अवमंदित} आयाम सिकुड़ते चरघातांकी \omterm{def:g10:refraction:fiber}{आवरण}
(बिंदुदार) के भीतर मरता जाता है।}
\end{omfigure}

\begin{definition}[प्रणोदित दोलन और अनुनाद]\label{def:g12:oscillators-and-time:resonance}
किसी \omterm{def:g12:oscillators-and-time:oscillator}{दोलित्र} को \omterm{def:g10:signals-and-waves:period}{आवर्ती} रूप से धकेलिए ---
\emph{प्रणोदित दोलन}% \index{प्रणोदित दोलन} --- और एक छोटे क्षणिक दौर के
बाद वह \emph{प्रणोदक} आवृत्ति $f$ पर दोलन करने लगता है, अपनी वाली
पर नहीं। उसका आयाम तब चरम पर होता है जब $f$ मुक्त \omterm{def:g12:oscillators-and-time:oscillator}{दोलित्र}
की \emph{स्वाभाविक आवृत्ति}% \index{स्वाभाविक आवृत्ति} $f_0$
के पास पहुँचे: यही \emph{अनुनाद}\index{अनुनाद} है, जिसमें हर धक्का गति के
क़दम से क़दम मिलाकर आता है और हर चक्र में ऊर्जा भरता जाता है; और
\omterm{def:g12:oscillators-and-time:damping}{अवमंदन} जितना हलका, शिखर उतना ही ऊँचा और तीखा।
\end{definition}

\begin{omfigure}
\begin{tikzpicture}
\begin{axis}[omaxis, width=8.6cm, height=4.6cm, xlabel={$f/f_0$},
    ylabel={आयाम (मनमाना)}, xmin=0, xmax=2.2, ymin=0, ymax=5.6,
    legend pos=north east]
  \addplot[omThm, thick, domain=0:2.2, samples=200] {1/sqrt((1-x^2)^2 + (x/5)^2)};
  \addplot[omDef, thick, domain=0:2.2, samples=120] {1/sqrt((1-x^2)^2 + (x/1.2)^2)};
  \legend{light damping, heavy damping}
  \addplot[omExo, dashed, forget plot] coordinates {(1,0) (1,5.4)};
\end{axis}
\end{tikzpicture}

{\small अनुनाद-वक्र: प्रणोदक आवृत्ति के सामने प्रणोदित
आयाम। $f = f_0$ के पास प्रतिक्रिया मीनार सी उठ जाती है, और
\omterm{def:g12:oscillators-and-time:damping}{अवमंदन} जितना हलका उतनी ही ऊँची।}
\end{omfigure}

\begin{example}[अनुनाद, भले के लिए और बुरे के लिए]\label{ex:g12:oscillators-and-time:resonance-zoo}
झूला झुलाता कोई अभिभावक उसे उसकी अपनी आवृत्ति पर धकेलता है --- यानी
काम पर लगा \omterm{def:g12:oscillators-and-time:resonance}{अनुनाद}। \num{2000} में लंदन का नया मिलेनियम पुल तब
ख़तरनाक ढंग से डोल उठा जब भीड़ के क़दम \qty{1}{Hz} के पास की किसी \omterm{def:g12:oscillators-and-time:resonance}{स्वाभाविक
आवृत्ति} से बँध गए; उसे अवमंदक लगाने के लिए दो साल बंद रखना पड़ा। और हर
क्वार्ट्ज़ घड़ी अपने क्रिस्टल को \qty{32768}{Hz} पर \omterm{def:g12:oscillators-and-time:resonance}{अनुनाद} में चलाकर गाता हुआ
बनाए रखती है।
\end{example}

\section{समय नापना}

\begin{method}[आवर्तकाल नापना]\label{met:g12:oscillators-and-time:timing}
एक दोलन का समय कभी मत नापिए: गोलक के \omterm{def:g10:universal-gravitation:weight}{ऊर्ध्वाधर} पार करते समय (जहाँ वह सबसे तेज़
है और परखना सबसे आसान) स्टॉपवॉच चालू कीजिए, $N = 50$ दोलन गिनिए, और $N$ से
भाग दीजिए --- \qty{0.2}{s} की प्रतिक्रिया-त्रुटि $T_0$ पर \qty{4}{ms} रह जाती है।
दोहराइए, और दौरों की आपस में तुलना कीजिए।
\end{method}

\begin{example}[टिक की तीन सदियाँ]\label{ex:g12:oscillators-and-time:clocks}
हाइगेंस ने पहली लोलक-घड़ी \num{1657} में बनाई: समकालिकता ने घड़ियों को दिन भर के
चौथाई घंटे के बहकाव से सेकंडों तक ले आया, और क्षतिपूरित लोलकों ने दिन भर में
\omterm{def:g10:orders-of-magnitude:unit}{सेकंड} के अंशों तक। क्वार्ट्ज़ की कलाई-घड़ी (\num{1969}) ने छड़ की जगह \qty{32768}{Hz} पर
लचकता क्रिस्टल रख दिया: महीने भर में \omterm{def:g10:orders-of-magnitude:unit}{सेकंड} के दसवें हिस्से। और \num{1967} से
\omterm{def:g10:orders-of-magnitude:unit}{सेकंड} ख़ुद एक परमाणविक दोलन से \emph{परिभाषित} होता है: सीज़ियम-133 के एक
आंतरिक संक्रमण के सूक्ष्मतरंग विकिरण के \num{9192631770} आवर्तकालों की अवधि
(\cref{ch:g12:quantum-world})। सबसे अच्छी सीज़ियम घड़ियाँ दस करोड़ वर्ष में एक \omterm{def:g10:orders-of-magnitude:unit}{सेकंड} बहकती हैं।
\end{example}

\begin{remark}[नैनोसेकंड के पीछे क्यों]\label{rem:g12:oscillators-and-time:gps}
GPS का ग्राही उपग्रहों से आते रेडियो संकेतों का समय नापकर
अपनी स्थिति खोजता है; और प्रकाश प्रति नैनोसेकंड \qty{30}{cm} चलता है, इसलिए एक
माइक्रोसेकंड की घड़ी-त्रुटि आपको \qty{300}{m} दूर बिठा देती है। इसीलिए
\omterm{def:g10:universal-gravitation:satellite}{उपग्रह} परमाणु-घड़ियाँ ढोते हैं --- और उस परिशुद्धता पर किसी घड़ी की चाल
इस पर निर्भर करने लगती है कि वह कितनी तेज़ चलती है और कितनी ऊँचाई पर बैठती है,
जिससे हम \cref{ch:g12:special-relativity} में मिलेंगे।
\end{remark}

\section{अभ्यास}

\begin{exercise}[$\star$]\label{exo:g12:oscillators-and-time:1}
(क) \qty{1.00}{m} के \omterm{def:g12:oscillators-and-time:pendulum}{सरल लोलक} का \omterm{def:g12:oscillators-and-time:oscillator}{आवर्तकाल} और आवृत्ति
निकालिए। (ख) लंबाई चौगुनी कर दी जाती है: \omterm{def:g12:oscillators-and-time:oscillator}{आवर्तकाल} का क्या होता है?
\end{exercise}

\begin{exercise}[$\star$]\label{exo:g12:oscillators-and-time:2}
विश्राम में हृदय प्रति मिनट $72$ बार धड़कता है; और व्याध-पतंग का पंख
\qty{30}{Hz} पर। हृदय की आवृत्ति और \omterm{def:g12:oscillators-and-time:oscillator}{आवर्तकाल} तथा पंख का
\omterm{def:g12:oscillators-and-time:oscillator}{आवर्तकाल} दीजिए; और अध्याय का कौन-सा शब्द प्रति धड़कन छाती के अधिकतम
उठाव का नाम है?
\end{exercise}

\begin{exercise}[$\star$]\label{exo:g12:oscillators-and-time:3}
\qty{0.10}{kg} का फिसलक \qty{25}{N/m} की स्प्रिंग पर सवार है: $T_0$ और $f_0$ निकालिए।
द्रव्यमान दुगुना करने पर \omterm{def:g12:oscillators-and-time:oscillator}{आवर्तकाल} का क्या होता है?
\end{exercise}

\begin{exercise}[$\star$]\label{exo:g12:oscillators-and-time:4}
किसी \omterm{def:g12:oscillators-and-time:damping}{अवमंदित} \omterm{def:g12:oscillators-and-time:oscillator}{दोलित्र} के क्रमागत उच्चिष्ठ: $t = 0$, $2.0$,
$4.0$, \qty{6.0}{s} पर $2.0$, $1.5$, $1.13$, \qty{0.85}{cm}।
\omterm{def:g12:oscillators-and-time:damping}{छद्म-आवर्तकाल} पढ़िए; दिखाइए कि हर \omterm{def:g12:oscillators-and-time:oscillator}{आवर्तकाल} में आयाम
एक नियत गुणक से गिरता है; और अगले उच्चिष्ठ का पूर्वानुमान लगाइए।
\end{exercise}

\begin{exercise}[$\star$]\label{exo:g12:oscillators-and-time:5}
कोई अंतरिक्ष यात्री एक लोलक और एक \omterm{prop:g12:oscillators-and-time:spring-period}{स्प्रिंग--द्रव्यमान दोलित्र} चंद्रमा
पर ले जाता है ($g = \qty{1.62}{m/s^2}$)। हर \omterm{def:g12:oscillators-and-time:oscillator}{आवर्तकाल} कितने गुना बदलता है? समझाइए।
\end{exercise}

\begin{exercise}[$\star\star$]\label{exo:g12:oscillators-and-time:6}
(क) जाँचिए कि $\sqrt{L/g}$ की विमा समय की है। (ख) $m$, $L$, $g$ से बना
$T_0$ का कोई सूत्र द्रव्यमान क्यों नहीं रख सकता? (ग) क्या यह तर्क कभी
$2\pi$ पैदा कर सकता है?
\end{exercise}

\begin{exercise}[$\star\star$]\label{exo:g12:oscillators-and-time:7}
दो बार अवकलन करके सत्यापित कीजिए कि $x(t) = A\cos(2\pi t/T_0 +
\varphi)$ $\frac{d^2x}{dt^2} = -(k/m)\,x$ को ठीक तभी संतुष्ट करता है
जब $T_0 = 2\pi\sqrt{m/k}$।
\end{exercise}

\begin{exercise}[$\star\star$]\label{exo:g12:oscillators-and-time:8}
\qty{40}{N/m} की \omterm{def:g12:oscillators-and-time:spring}{दृढ़ता} वाली स्प्रिंग \qty{0.20}{kg} के फिसलक को \qty{3.0}{cm}
आयाम के साथ ढोती है। \omterm{def:g11:mechanical-energy:em}{यांत्रिक ऊर्जा} और अधिकतम चाल निकालिए।
चाल कहाँ अधिकतम है? कहाँ शून्य?
\end{exercise}

\begin{exercise}[$\star\star$]\label{exo:g12:oscillators-and-time:9}
(क) \omterm{def:g10:orders-of-magnitude:unit}{सेकंड} वाले लोलक की लंबाई निकालिए ($T_0 =
\qty{2.00}{s}$)। (ख) उसकी लंबाई $1.0\%$ बढ़
जाती है: $\sqrt{1+x}
\approx 1 + x/2$ के साथ $T_0$ का सापेक्ष बदलाव और दैनिक त्रुटि निकालिए।
\end{exercise}

\begin{exercise}[$\star\star$]\label{exo:g12:oscillators-and-time:10}
कोई हलका \omterm{def:g12:oscillators-and-time:damping}{अवमंदित} \omterm{def:g12:oscillators-and-time:oscillator}{दोलित्र} हर \omterm{def:g12:oscillators-and-time:oscillator}{आवर्तकाल} में अपनी
ऊर्जा का $5.0\%$ गँवा देता है। (क) प्रति \omterm{def:g12:oscillators-and-time:oscillator}{आवर्तकाल}
आयाम की गिरावट? (ख) आधी ऊर्जा गँवाने में कितने
\omterm{def:g12:oscillators-and-time:oscillator}{आवर्तकाल} लगेंगे? (ग) क्या \omterm{def:g12:oscillators-and-time:damping}{छद्म-आवर्तकाल} $T_0$ से बहुत दूर
है?
\end{exercise}

\begin{exercise}[$\star\star$]\label{exo:g12:oscillators-and-time:11}
अनुनाद-वक्र से समझाइए: (क) झूला झुलाना उसकी अपनी लय पर ही क्यों
काम करता है; (ख) सिपाही पैदल-पुल पर क़दम-ताल क्यों तोड़ देते हैं;
(ग) कपड़े धोने की मशीन घूमना बढ़ाते समय एक ही चाल पर क्यों काँप उठती है।
\end{exercise}

\begin{exercise}[$\star\star\star$]\label{exo:g12:oscillators-and-time:12}
समुद्र-तल पर ठीक चलती एक लोलक-घड़ी ऐसी वेधशाला में ले जाई जाती है जहाँ $g = \qty{9.78}{m/s^2}$।
(क) तेज़ या धीमी? प्रतिदिन कितने \omterm{def:g10:orders-of-magnitude:unit}{सेकंड}? (ख) उसके \qty{0.994}{m} के लोलक को कितना छोटा
करना पड़ेगा?
\end{exercise}

\begin{exercise}[$\star\star\star$]\label{exo:g12:oscillators-and-time:13}
\qty{0.20}{kg} का फिसलक $x(t) = A\cos(2\pi t/T_0)$ मानता है, जहाँ $A =
\qty{4.0}{cm}$, $T_0 = \qty{0.63}{s}$। अधिकतम चाल, अधिकतम
त्वरण, \omterm{def:g12:oscillators-and-time:spring}{दृढ़ता} $k$ और \omterm{def:g11:mechanical-energy:em}{यांत्रिक ऊर्जा} निकालिए।
\end{exercise}

\begin{exercise}[$\star\star\star$]\label{exo:g12:oscillators-and-time:14}
$g$ नापने के लिए: $L = 99.5 \pm \qty{0.2}{cm}$ वाला लोलक $40.1 \pm \qty{0.2}{s}$ में $20$ दोलन करता है।
$g$ और उसकी \omterm{def:g10:orders-of-magnitude:uncertainty}{सापेक्ष अनिश्चितता} निकालिए ($L$ और $T_0^2$ की
अनिश्चितताएँ जोड़ लीजिए); क्या परिणाम \qty{9.81}{m/s^2} से मेल खाता है?
\end{exercise}

\begin{exercise}[$\star\star\star$]\label{exo:g12:oscillators-and-time:15}
समय रखने वालों की तुलना सापेक्ष बहकाव (बीते समय पर त्रुटि) से कीजिए: प्रतिदिन
\qty{10}{s} गँवाती लोलक-घड़ी, महीने भर में \qty{0.5}{s} वाली क्वार्ट्ज़ घड़ी, और दस करोड़
वर्ष में \qty{1}{s} वाली सीज़ियम घड़ी: तीनों अनुपात निकालिए। \qty{1}{\micro s} से ग़लत GPS
घड़ी: स्थिति में त्रुटि? GPS को किस कुल की ज़रूरत है?
\end{exercise}

\section{समस्या: घड़ीसाज़ की मेज़}

\begin{problem}[{सप्ताहांत समस्या --- घड़ीसाज़ की मेज़: दीवार-घड़ी को फिर से सही
समय पर लाना --- पचास झूलों का समय, गरमियों में फैलती पीतल की छड़, निर्गम-तंत्र
का फुसफुसाता धक्का, और क्वार्ट्ज़ का एक चुनौती-देने वाला --- और फ़ैसला, प्रति
महीना सेकंडों में}]
\label{pb:g12:oscillators-and-time:1}
कोई लोलक वाली दीवार-घड़ी ``ग़लत चलती'' मेज़ पर आती है। उसकी पीतल की छड़ और भारी
गोलक को प्रभावी लंबाई $L$ और गोलक-द्रव्यमान $m = \qty{1.2}{kg}$ वाला \omterm{def:g12:oscillators-and-time:pendulum}{सरल
लोलक} मानिए; और दाँतेदार पहिए हर पूरे दोलन को ठीक \qty{2.000}{s} गिनते हैं।

\medskip
\noindent\textbf{भाग I --- घड़ी की नब्ज़ देखना।}
\begin{enumerate}
  \item एक के बजाय $50$ दोलनों का समय क्यों नापा जाए? यदि आपका
        प्रतिक्रिया-समय \qty{0.2}{s} हो तो $T_0$ पर बची त्रुटि का अनुमान
        लगाइए।
  \item स्टॉपवॉच $50$ दोलनों के लिए \qty{99.4}{s} पढ़ती है: \omterm{def:g12:oscillators-and-time:oscillator}{आवर्तकाल}?
  \item लोलक की वर्तमान प्रभावी लंबाई निकालिए।
  \item कौन-सी लंबाई ठीक \qty{2.000}{s} देती? \omterm{def:g11:lenses-and-eye:eye}{समंजन} का पेंच गोलक को ऊपर उठाए या नीचे
        करे, और कितना?
  \item बिना \omterm{def:g11:lenses-and-eye:eye}{समंजन} घड़ी तेज़ है या धीमी, और प्रतिदिन कितने \omterm{def:g10:orders-of-magnitude:unit}{सेकंड}?
  \item $g$ के प्रति संवेदनशीलता: चंद्रमा पर ($g = \qty{1.62}{m/s^2}$)
        \omterm{def:g12:oscillators-and-time:oscillator}{आवर्तकाल}? और उस वेधशाला में दैनिक त्रुटि जहाँ $g = \qty{9.80}{m/s^2}$?
\end{enumerate}

\medskip
\noindent\textbf{भाग II --- गरमियों की ढील।} पीतल फैलता है: $\Delta\theta$ गरम होने
पर छड़ खिंचकर $L' = L\,(1 + \lambda\,
\Delta\theta)$ हो जाती है, प्रति \omterm{def:g10:light-spectra:kelvin}{केल्विन} $\lambda = \num{1.9e-5}$। घड़ी \qty{20}{\celsius}
पर समंजित की गई है।
\begin{enumerate}[resume]
  \item \qty{25}{\celsius} वाले गरमियों के कमरे के लिए $\Delta L$ निकालिए।
  \item $\sqrt{1+x} \approx 1 + x/2$ का उपयोग करके दिखाइए कि $\Delta T_0 / T_0 = \tfrac12 \lambda\, \Delta\theta$।
  \item $\Delta T_0/T_0$ और गरमियों की दैनिक त्रुटि निकालिए: तेज़ या धीमी?
  \item \qty{10}{\celsius} वाली सर्दियों की गैलरी के लिए वही प्रश्न।
  \item पुरानी परिशुद्ध घड़ियाँ ``ग्रिडआयरन'' लोलक इस्तेमाल करती थीं, जिनमें
        पीतल की छड़ें इस्पात की छड़ों के साथ (जो लगभग आधा फैलती हैं) उलटी-सीधी
        दिशाओं में लगी होती थीं। समझाइए कि इससे चाल कैसे थमी रह सकती है।
\end{enumerate}

\medskip
\noindent\textbf{भाग III --- फुसफुसाता धक्का।} निर्गम-तंत्र अलग कर देने पर
आयाम $\theta_m = 2.00^\circ$ से शुरू किया गया झूला \qty{280}{s} में अपना आयाम
आधा कर लेता है; और सेवा में, हर \omterm{def:g12:oscillators-and-time:oscillator}{आवर्तकाल} पर निर्गम-तंत्र का ज़रा-सा
धक्का $\theta_m$ को अचर बनाए रखता है।
\begin{enumerate}[resume]
  \item आधा होने में कितने \omterm{def:g12:oscillators-and-time:oscillator}{आवर्तकाल} लगते हैं?
  \item छोर पर गोलक $h = L(1 - \cos\theta_m)
        \approx L\,\theta_m^2/2$ (रेडियन) चढ़ चुका होता है: इससे $E_m = \tfrac12
        m g L\,\theta_m^2$ निकालिए
        और उसका मान लीजिए (\cref{ch:g11:mechanical-energy})।
  \item ऊर्जा आयाम के वर्ग के समानुपाती है: प्रति
        \omterm{def:g12:oscillators-and-time:oscillator}{आवर्तकाल} $E_m$ का कितना अंश गँवाया जाता है?
  \item हर धक्के की ऊर्जा और ज़रूरी औसत शक्ति निकालिए।
  \item चालक भार $M = \qty{2.5}{kg}$ हर सप्ताह $h =
        \qty{1.0}{m}$ नीचे उतरता है: क्या
        इतनी ऊर्जा पर्याप्त है, और किस \omterm{def:g10:energy-conservation:efficiency}{दक्षता} पर?
\end{enumerate}

\medskip
\noindent\textbf{भाग IV --- क्वार्ट्ज़ की चुनौती।} कोई क्वार्ट्ज़ यंत्र $f = \qty{32768}{Hz}$
पर दोलन करता है और महीने भर में लगभग \qty{0.3}{s} बहकता है; और सुधरी हुई लोलक-घड़ी
अब भी प्रतिदिन \qty{0.5}{s} बहकती है।
\begin{enumerate}[resume]
  \item क्वार्ट्ज़ का \omterm{def:g12:oscillators-and-time:oscillator}{आवर्तकाल} निकालिए। दिखाइए कि $\num{32768} = 2^{15}$: घड़ी के
        परिपथ को दो की घात ही क्यों चाहिए?
  \item क्वार्ट्ज़ और सुधरी लोलक-घड़ी के सापेक्ष बहकाव निकालिए, और दोनों को
        प्रति महीना सेकंडों में भी। कौन जीतता है, और कितने गुना से?
  \item सीज़ियम मानक --- जिसके \num{9192631770} \omterm{def:g12:oscillators-and-time:oscillator}{आवर्तकाल} एक \omterm{def:g10:orders-of-magnitude:unit}{सेकंड} बनाते हैं ---
        दस करोड़ वर्ष में \qty{1}{s} बहकता है: उसका सापेक्ष बहकाव? वह क्वार्ट्ज़
        से कितनी कोटि नीचे है, और GPS उसी की माँग क्यों करता है?
  \item फ़ैसला, संख्याओं समेत एक वाक्य में: सुधरी हुई घड़ी दीवार पर क्या कमा
        लेती है, और \omterm{def:g10:orders-of-magnitude:unit}{सेकंड} का मालिक अब कौन है?
\end{enumerate}
\end{problem}
